\documentclass[10pt,a4paper]{amsart}
\usepackage[margin=19mm]{geometry}
\usepackage{amsmath,amssymb,mathtools}
\usepackage[scr=rsfs,cal=euler]{mathalfa}
\usepackage{xfrac}
\usepackage[unicode,hidelinks,bookmarks=false]{hyperref}
\newcommand{\ie}{i.e.\ }
\newcommand{\cf}{cf.\ }
\newcommand{\resp}{respectively\ }
\newcommand{\Subsection}{\subsection}
\providecommand{\nobreakdash}{\nobreak\hbox{-}}
\setlength{\emergencystretch}{2em}
\multlinegap0pt
\usepackage{bm}
\usepackage{bbm}
\usepackage{dsfont}
\usepackage{xspace}
\usepackage{cancel}
\usepackage{mathtools}
\usepackage{cleveref}
\usepackage{amsthm}
\makeatletter
\@ifundefined{definition}{\newtheorem{definition}{Definition}}{}
\@ifundefined{lemma}{\newtheorem{lemma}{Lemma}}{}
\@ifundefined{proposition}{\newtheorem{proposition}{Proposition}}{}
\makeatother
\newcommand{\KK}{\mathcal{K}}
\newcommand{\LL}{\mathcal{L}}
\title[Cheeger Inequalities for the Persistent Laplacian]{Cheeger Inequalities for the Persistent Laplacian}
\author{Magnus Bakke Botnan, Rui Dong}
\date{Source: arXiv:2606.02846v2 (CC BY 4.0)}
\begin{document}
\maketitle

\noindent\emph{Source and licence.} Cheeger Inequalities for the Persistent Laplacian. Version arXiv:2606.02846v2 (CC BY 4.0), distributed under the Creative Commons Attribution 4.0 International licence, as recorded in the arXiv licence metadata for this exact version and shown on the abstract page https://arxiv.org/abs/2606.02846v2. Full rights notice: licenses/arxiv-2606.02846-CC-BY-4.0.txt.\par\smallskip

\noindent\textit{Editorial scope.} Counted unit: the Proposition whose source label is \texttt{prop\_cheeger\_ineq} in the article (source0.tex lines 1014--1027, source label \texttt{prop\_cheeger\_ineq}), delivered together with the article's own complete original proof of it (source0.tex lines 1029--1098, one \texttt{proof} environment of 70 lines). No mathematics is written, repaired or re-derived: statement and proof are the article's own text line for line. Required context retained uncounted: lem:p1cheeger\_equiv (lines 894--919); def:nonzero\_cheeger (lines 982--991). Every definition, display and label of those lines is retained unchanged. Omitted from this package and not redistributed: 1--893, 920--981, 992--1013, 1028--1028, 1099--4167. Nothing inside a retained excerpt is omitted or altered: every retained body is delivered exactly as the article's lines are written, so no omission receipt of any kind accompanies this package. Citations: the retained excerpts carry no citation command at all, so no bibliography entry of the article is transcribed and no reference is redistributed. Reference adaptation: none was needed and none was made; every reference inside the delivered unit resolves inside it, so no unresolved reference or citation is produced. Printed slips kept literal: no printed word, symbol, hypothesis or number of the counted statement or of its proof was altered. Layout and class: the article's own class is \texttt{article} with packages including babel, geometry, amsmath, mathtools, graphicx, amssymb, hyperref, amsthm, enumitem, tikz, tikz-cd, todonotes, pgfplots, adjustbox; the wrapper is the lane's pinned \texttt{amsart} editorial wrapper at 10pt on A4, which restates the theorem-like environments the retained text needs and adds the article's own macro definitions that the retained text needs (\texttt{\textbackslash KK}, \texttt{\textbackslash LL}), with extra wrapper packages (bm, bbm, dsfont, xspace, cancel, mathtools, cleveref). The delivered numbering is the unit's own. Assets: no retained span contains a figure environment, a graphics-inclusion command, a PDF-page inclusion, a table or a TikZ picture; no image file is copied into this package, the package declares no asset dependency and no image file is redistributed with it. Licence: version arXiv:2606.02846v2 of this article is distributed under the Creative Commons Attribution 4.0 International licence, as recorded in the arXiv licence metadata for this exact version and shown on the abstract page https://arxiv.org/abs/2606.02846v2. A full rights notice is delivered at licenses/arxiv-2606.02846-CC-BY-4.0.txt. Characters: every retained excerpt is pure ASCII, so the delivered engine, XeLaTeX with its default Latin Modern text font, reports no missing character.
\begin{lemma}
\label{lem:p1cheeger_equiv}
Let
\[
i=\dim\ker(PB^T)+1.
\]
Then
\[
\lambda_i\bigl(\Delta_{q,1,\mathrm{up}}^{\KK,\LL}\bigr)
=
\inf_{\gamma(S)\geq i}
\sup_{x\in S}
\frac{\|PB^T x\|_1}{\|x\|_1}
=
\min_{\substack{x\neq 0\\ x\perp^1\ker(PB^T)}}
\frac{\|PB^T x\|_1}{\|x\|_1}.
\]
Here \(x\perp^1 Y\) means that
\[
\|x+y\|_1\geq \|x\|_1
\qquad
\text{for every } y\in Y.
\]
Equivalently, \(x\) is an \(\ell^1\)-minimal representative of the affine
space \(x+Y\).
\end{lemma}

\begin{definition}\label{def:nonzero_cheeger}
Let \(\KK\hookrightarrow \LL\) be an inclusion of simplicial complexes. The
persistent Cheeger constant of this inclusion is
\[
\varphi_q^{\KK,\LL}
=
\min_{\substack{x\neq 0\\ x\perp^1\ker(PB^T)}}
\frac{\|PB^T x\|_1}{\|x\|_1}.
\]
\end{definition}

\begin{proposition}\label{prop_cheeger_ineq}
Let \(\KK\hookrightarrow \LL\) be an inclusion of simplicial complexes, and
let \(|S_q^\KK|\) and \(|S_{q+1}^\LL|\) denote the numbers of \(q\)-simplices
in \(\KK\) and \((q+1)\)-simplices in \(\LL\), respectively. Let
\(\lambda_{\min}^+\) denote the smallest non-zero eigenvalue of
the persistent up Laplacian \(\Delta_{q,\mathrm{up}}^{\KK,\LL}\). Then
\begin{equation}\label{ineq:persistent_cheeger}
\frac{\bigl(\varphi_q^{\KK,\LL}\bigr)^2}{|S_{q+1}^{\LL}|}
\leq
\lambda_{\min}^+
\leq
|S_q^\KK|\bigl(\varphi_q^{\KK,\LL}\bigr)^2 .
\end{equation}
\end{proposition}

\begin{proof}
Let \(x\in C_q^\KK\) be nonzero. We use the elementary norm inequalities
\[
1\leq \frac{\|x\|_1^2}{\|x\|_2^2}\leq |S_q^\KK|,
\qquad
1\leq
\frac{\|PB^T x\|_1^2}{\|PB^T x\|_2^2}
\leq |S_{q+1}^{\LL}|,
\]
whenever \(PB^T x\neq 0\). Hence
\[
\frac{1}{|S_{q+1}^{\LL}|}
\frac{\|PB^T x\|_1^2}{\|x\|_1^2}
\leq
\frac{\|PB^T x\|_2^2}{\|x\|_2^2}
\leq
|S_q^\KK|
\frac{\|PB^T x\|_1^2}{\|x\|_1^2}.
\]

Set
\[
i=\dim\ker(PB^T)+1.
\]
By the discussion above, this is precisely the index of the smallest non-zero variational eigenvalue.By \cref{lem:p1cheeger_equiv} and \cref{def:nonzero_cheeger},
\[
\varphi_q^{\KK,\LL}
=
\lambda_i\bigl(\Delta_{q,1,\mathrm{up}}^{\KK,\LL}\bigr)
=
\inf_{\gamma(S)\geq i}
\sup_{x\in S}
\frac{\|PB^T x\|_1}{\|x\|_1}.
\]
On the other hand,
\[
\lambda_i\bigl(\Delta_{q,\mathrm{up}}^{\KK,\LL}\bigr)
=
\inf_{\gamma(S)\geq i}
\sup_{x\in S}
\frac{\|PB^T x\|_2^2}{\|x\|_2^2}.
\]
Applying the preceding norm comparison to the Rayleigh quotients gives
\[
\frac{1}{|S_{q+1}^{\LL}|}
\left(
\inf_{\gamma(S)\geq i}
\sup_{x\in S}
\frac{\|PB^T x\|_1}{\|x\|_1}
\right)^2
\leq
\lambda_i\bigl(\Delta_{q,\mathrm{up}}^{\KK,\LL}\bigr)
\leq
|S_q^\KK|
\left(
\inf_{\gamma(S)\geq i}
\sup_{x\in S}
\frac{\|PB^T x\|_1}{\|x\|_1}
\right)^2.
\]
Using the identity above with \(\varphi_q^{\KK,\LL}\), we obtain
\[
\frac{\bigl(\varphi_q^{\KK,\LL}\bigr)^2}{|S_{q+1}^{\LL}|}
\leq
\lambda_i\bigl(\Delta_{q,\mathrm{up}}^{\KK,\LL}\bigr)
\leq
|S_q^\KK|\bigl(\varphi_q^{\KK,\LL}\bigr)^2,
\]
as claimed.
\end{proof}

\end{document}
